\section{Conclusion}
\label{sec:conclusion}
This paper investigated whether alternative sentiment data contains incremental predictive power for forecasting realized volatility beyond the established HAR framework. Our empirical results suggest that directly augmenting linear models with high-dimensional sentiment scores fails out-of-sample. However, when these features are processed through a non-linear deep learning architecture (Neural-HAR) comprising a Gated Residual Network and an FT-Transformer, we observe a statistically significant improvement in forecasting accuracy, primarily under the QLIKE loss function. 

Despite this statistical edge and its ability to act as the sole survivor in the S\&P 500 Model Confidence Set, the economic benefits of the Neural-HAR model are nuanced. The model does not generate a higher Sharpe ratio due to the transaction costs associated with its elevated turnover. Instead, its economic value lies in superior risk mitigation, specifically by yielding tighter realized volatility profiles and shallower maximum drawdowns in a volatility-targeting framework. Furthermore, we find minimal evidence supporting a strictly regime-conditional sentiment effect. Ultimately, while non-linear architectures can successfully extract volatility signals from alternative data, the practical translation of this signal into net profitability remains highly constrained by market frictions.
